\documentclass{amsart}
% For dealing with references we use the comment environment
\usepackage{verbatim}
\newenvironment{reference}{\comment}{\endcomment}
%\newenvironment{reference}{}{}
\newenvironment{slogan}{\comment}{\endcomment}
\newenvironment{history}{\comment}{\endcomment}

% For commutative diagrams we use Xy-pic
\usepackage[all]{xy}

% We use 2cell for 2-commutative diagrams.
\xyoption{2cell}
\UseAllTwocells

% We use multicol for the list of chapters between chapters
\usepackage{multicol}

% This is generally recommended for better output
\usepackage{lmodern}
\usepackage[T1]{fontenc}

% For cross-file-references

% Package for hypertext links:
\usepackage{hyperref}

% For any local file, say "hello.tex" you want to link to please

% Theorem environments.
%
\theoremstyle{plain}
\newtheorem{theorem}[subsection]{Theorem}
\newtheorem{proposition}[subsection]{Proposition}
\newtheorem{lemma}[subsection]{Lemma}

\theoremstyle{definition}
\newtheorem{definition}[subsection]{Definition}
\newtheorem{example}[subsection]{Example}
\newtheorem{exercise}[subsection]{Exercise}
\newtheorem{situation}[subsection]{Situation}

\theoremstyle{remark}
\newtheorem{remark}[subsection]{Remark}
\newtheorem{remarks}[subsection]{Remarks}

\numberwithin{equation}{subsection}

% Macros
%
\def\lim{\mathop{\mathrm{lim}}\nolimits}
\def\colim{\mathop{\mathrm{colim}}\nolimits}
\def\Spec{\mathop{\mathrm{Spec}}}
\def\Hom{\mathop{\mathrm{Hom}}\nolimits}
\def\Ext{\mathop{\mathrm{Ext}}\nolimits}
\def\SheafHom{\mathop{\mathcal{H}\!\mathit{om}}\nolimits}
\def\SheafExt{\mathop{\mathcal{E}\!\mathit{xt}}\nolimits}
\def\Sch{\mathit{Sch}}
\def\Mor{\mathop{\mathrm{Mor}}\nolimits}
\def\Ob{\mathop{\mathrm{Ob}}\nolimits}
\def\Sh{\mathop{\mathit{Sh}}\nolimits}
\def\NL{\mathop{N\!L}\nolimits}
\def\CH{\mathop{\mathrm{CH}}\nolimits}
\def\proetale{{pro\text{-}\acute{e}tale}}
\def\etale{{\acute{e}tale}}
\def\QCoh{\mathit{QCoh}}
\def\Ker{\mathop{\mathrm{Ker}}}
\def\Im{\mathop{\mathrm{Im}}}
\def\Coker{\mathop{\mathrm{Coker}}}
\def\Coim{\mathop{\mathrm{Coim}}}

% Boxtimes
%
\DeclareMathSymbol{\boxtimes}{\mathbin}{AMSa}{"02}

%
% Macros for moduli stacks/spaces
%
\def\QCohstack{\mathcal{QC}\!\mathit{oh}}
\def\Cohstack{\mathcal{C}\!\mathit{oh}}
\def\Spacesstack{\mathcal{S}\!\mathit{paces}}
\def\Quotfunctor{\mathrm{Quot}}
\def\Hilbfunctor{\mathrm{Hilb}}
\def\Curvesstack{\mathcal{C}\!\mathit{urves}}
\def\Polarizedstack{\mathcal{P}\!\mathit{olarized}}
\def\Complexesstack{\mathcal{C}\!\mathit{omplexes}}
% \Pic is the operator that assigns to X its picard group, usage \Pic(X)
% \Picardstack_{X/B} denotes the Picard stack of X over B
% \Picardfunctor_{X/B} denotes the Picard functor of X over B
\def\Pic{\mathop{\mathrm{Pic}}\nolimits}
\def\Picardstack{\mathcal{P}\!\mathit{ic}}
\def\Picardfunctor{\mathrm{Pic}}
\def\Deformationcategory{\mathcal{D}\!\mathit{ef}}

\setlength{\emergencystretch}{3em}
\begin{document}
\title[Stacks Project, Tag 0A7E]{244 --- The image of the unit under an A-linear coherent derived equivalence is invertible}
\maketitle
\noindent Copyright (C) 2005--2025 Johan de Jong. Stacks Project authors. Source: \href{https://stacks.math.columbia.edu/tag/0A7E}{Tag 0A7E}.

\noindent Editorial difficulty note (not author text). Difficulty H2: The argument first forces the image of every residue field to be concentrated in one degree by ruling out negative self-extensions. It then detects all remaining cohomology of the image of A using finite-length test modules, Nakayama, and endomorphism rings. Recovering local shifts of rank-one free modules from categorical data is a substantive reconstruction argument..

\noindent Editorial provenance note (not author text). History: excerpt from revision \texttt{a0444\allowbreak 6e57e\allowbreak c1fbc\allowbreak 252a8\allowbreak 71afc\allowbreak ec775\allowbreak 2fb28\allowbreak 07b14}, dualizing.tex, lines 2641--2699. Reference presentation was adapted; original-author context and core support blocks were retained or added. The mathematical wording is unchanged. Licensed under GNU FDL 1.2 or later; no invariant sections or cover texts. See ../licenses/COPYING. Context blocks reproduce author definitions and supporting results; they are not additional counted problems.

\begin{lemma}
\label{lemma-equivalence-comes-from-invertible}
Let $A$ be a Noetherian ring. Let
$F : D^b_{\textit{Coh}}(A) \to D^b_{\textit{Coh}}(A)$ be an $A$-linear
equivalence of categories. Then $F(A)$ is an invertible object of $D(A)$.
\end{lemma}

\begin{proof}
Let $\mathfrak m \subset A$ be a maximal ideal with residue field $\kappa$.
Consider the object $F(\kappa)$. Since
$\kappa = \Hom_{D(A)}(\kappa, \kappa)$ we find that all
cohomology groups of $F(\kappa)$ are annihilated by $\mathfrak m$.
We also see that
$$
\Ext^i_A(\kappa, \kappa) = \text{Ext}^i_A(F(\kappa), F(\kappa))
= \Hom_{D(A)}(F(\kappa), F(\kappa)[i])
$$
is zero for $i < 0$. Say $H^a(F(\kappa)) \not = 0$ and
$H^b(F(\kappa)) \not = 0$ with $a$ minimal and $b$ maximal
(so in particular $a \leq b$). Then there is a nonzero map
$$
F(\kappa) \to H^b(F(\kappa))[-b] \to H^a(F(\kappa))[-b]
\to F(\kappa)[a - b]
$$
in $D(A)$ (nonzero because it induces a nonzero map on cohomology).
This proves that $b = a$. We conclude that $F(\kappa) = \kappa[-a]$.

\medskip\noindent
Let $G$ be a quasi-inverse to our functor $F$. Arguing as above
we find an integer $b$ such that $G(\kappa) = \kappa[-b]$.
On composing we find $a + b = 0$. Let $E$ be a finite $A$-module
which is annihilated by a power of $\mathfrak m$. Arguing by
induction on the length of $E$ we find that $G(E) = E'[-b]$
for some finite $A$-module $E'$ annihilated by a power of
$\mathfrak m$. Then $E[-a] = F(E')$.
Next, we consider the groups
$$
\Ext^i_A(A, E') = \text{Ext}^i_A(F(A), F(E')) =
\Hom_{D(A)}(F(A), E[-a + i])
$$
The left hand side is nonzero if and only if $i = 0$ and then
we get $E'$. Applying this with $E = E' = \kappa$ and using Nakayama's
lemma this implies that $H^j(F(A))_\mathfrak m$ is zero for $j > a$ and
generated by $1$ element for $j = a$. On the other hand, if
$H^j(F(A))_\mathfrak m$ is not zero for some $j < a$, then
there is a map $F(A) \to E[-a + i]$ for some $i < 0$ and some
$E$ (More on Algebra, Lemma \href{https://stacks.math.columbia.edu/tag/0A7D}{lemma detect cohomology (Tag 0A7D)})
which is a contradiction.
Thus we see that $F(A)_\mathfrak m = M[-a]$
for some $A_\mathfrak m$-module $M$ generated by $1$ element.
However, since
$$
A_\mathfrak m = \Hom_{D(A)}(A, A)_\mathfrak m =
\Hom_{D(A)}(F(A), F(A))_\mathfrak m = \Hom_{A_\mathfrak m}(M, M)
$$
we see that $M \cong A_\mathfrak m$. We conclude that there exists
an element $f \in A$, $f \not \in \mathfrak m$ such that
$F(A)_f$ is isomorphic to $A_f[-a]$. This finishes the proof.
\end{proof}
\end{document}
